%%%%%%%%%%%%%%%%%%%%%%%%%%%%%%%%%%%%%%%%%%%%%%%%%%%%%%%%%%%%%%%%%%%%%%%%%%%%%%%%
% Top Achievements & Elevator Pitch — Dr. Dipankar Bhattacharya
% Compiled: achievements-highlights.pdf
%%%%%%%%%%%%%%%%%%%%%%%%%%%%%%%%%%%%%%%%%%%%%%%%%%%%%%%%%%%%%%%%%%%%%%%%%%%%%%%%
\documentclass[11pt,a4paper]{article}

\usepackage[T1]{fontenc}
\usepackage[utf8]{inputenc}
\usepackage[p,osf,swashQ]{cochineal}
\usepackage[medium,bold]{cabin}
\usepackage[english]{babel}
\usepackage[a4paper,margin=2.2cm,top=2.2cm,bottom=2.2cm]{geometry}
\usepackage{xcolor}
\usepackage{enumitem}
\usepackage{titlesec}
\usepackage{hyperref}

\definecolor{accent}{HTML}{1B365D}
\definecolor{charcoal}{HTML}{222222}
\definecolor{rankgold}{HTML}{B8860B}

\hypersetup{colorlinks=true,allcolors=accent,breaklinks=true}
\setlength{\parindent}{0pt}
\setlength{\parskip}{0.45em}
\pagestyle{empty}

\titleformat{\section}{\Large\bfseries\sffamily\color{accent}}{}{0pt}{}[\vspace{-0.2em}\rule{\linewidth}{0.8pt}]
\titlespacing*{\section}{0pt}{1.1em}{0.45em}

\newcommand{\rank}[1]{\textcolor{rankgold}{\textbf{#1}}}
\newcommand{\btit}[1]{\textbf{\textit{#1}}}

\begin{document}

\begin{center}
  {\Huge\bfseries\sffamily\color{accent} Top Achievements}\\[0.35em]
  {\large\bfseries Dr. Dipankar Bhattacharya, Ph.D.}\\[0.2em]
  {\normalsize Marie Sk\l{}odowska-Curie Postdoctoral Fellow, Imperial College London}\\[0.15em]
  {\small Dyson School of Design Engineering \textbar{} March 2026}
\end{center}

\vspace{0.4em}

\section*{Purpose}
This note distils the strongest evidence from my CV for promotion, fellowship, and panel discussions.
Items are ranked by \textbf{impact}: funding scale, venue prestige, independence, translation to real-world use, and external recognition.

\section*{Overall ranking --- top 14}

\begin{enumerate}[leftmargin=1.6em, itemsep=0.55em, label=\rank{\arabic*.}]

  \item \textbf{Horizon Europe Marie Sk\l{}odowska-Curie Postdoctoral Fellowship (CASREx).}
  Grant No.~101205678; EUR~276,187.92; 24~months; Imperial College London.
  I lead an independent fellowship on cable-driven exosuits for upper-arm muscle-impairment rehabilitation.

  \item \textbf{Research independence across three domains.}
  Soft medical robotics (PhD, Auckland) $\to$ cable-driven robots (CUHK) $\to$ garment robotics with imitation learning (TransGP) $\to$ rehabilitation exosuits (Imperial CASREx).

  \item \textbf{Industry-funded translation --- BvW Holding stent testing on RoSE.}
  NZD~39,000 industry contract (2018); outputs in \textit{Soft Robotics} (2021) and \textit{IEEE TIE} (2021).

  \item \textbf{Flagship journal pipeline (first-author core).}
  \textit{Mechanism and Machine Theory} (2025); \textit{IEEE TCST} (2023); \textit{Soft Robotics} RoSE (2021); \textit{IEEE TIE} (2020, 2021).

  \item \textbf{Recent high-profile acceptances in manufacturing robotics.}
  \textit{IEEE T-ASE} (2026, fabric alignment); \textit{IEEE IROS} (2026, RTFF, equal contribution); \textit{IEEE RA-L} (2025, equal contribution).

  \item \textbf{Led multidisciplinary garment-robotics programme at TransGP} (Jul 2024--Oct 2025).
  Dual-arm fabric testbeds; mesh perception + imitation learning; contributions to sewing, flattening, and destacking workflows.

  \item \textbf{European Talent Academy 2026} (Imperial College London, TU Munich, Politecnico di Milano).
  Selected Dec 2025; Munich session completed Jun 2026.

  \item \textbf{Invited speaker, World Robotics Conference 2025, Beijing.}
  Talk: \textit{Fabric modeling and perception for garment manufacturing}.

  \item \textbf{PhD thesis output density} (Auckland, 2016--2021).
  RoSE platform; 3 journal articles, 2 conference papers, 1 Springer book chapter (2019).

  \item \textbf{Postdoctoral research track record at CUHK} (2021--2024).
  MMT and TCST papers; grant co-authorship; mentoring of UG and PhD researchers.

  \item \textbf{Peer review service} (2021--present).
  Reviewer for \textit{IEEE RA-L}, \textit{ICRA}, and \textit{ASME JMR}.

  \item \textbf{Best Conference Paper Award, RiTA 2017}, Daejeon, Korea.

  \item \textbf{GATE 99.12\% percentile} (Rank 1,211 / 137,853), 2011.

  \item \textbf{Three years lecturing} (Assistant Professor Grade-1), Galgotias University, 2013--2016.

\end{enumerate}

\section*{By framework category}

\subsection*{\sffamily\color{accent} Funding \& grants}
\begin{itemize}[leftmargin=1.4em, itemsep=0.35em]
  \item MSCA Postdoctoral Fellowship --- \textbf{EUR~276,187.92} (lead)
  \item BvW Holding RoSE stent testing --- \textbf{NZD~39,000} (lead)
  \item ITF PRP and RTH-ITF grants --- co-author (CUHK)
  \item PhD scholarships --- University of Auckland; Riddet CoRE; MHRD GATE (IIT)
\end{itemize}

\subsection*{\sffamily\color{accent} Research outputs (quality \& quantity)}
\textbf{Journals:} IEEE T-ASE (2026, accepted); IEEE RA-L (2025); MMT (2025); IEEE TCST (2023); Soft Robotics (2021, RoSE \& SoGut); IEEE TIE (2020, 2021).\\
\textbf{Conferences:} IROS 2026 (accepted); IEEE ICMA 2025; ICRA 2023 poster; M2VIP 2017/2018; Springer book chapter (2019).

\subsection*{\sffamily\color{accent} Research independence \& leadership}
\begin{itemize}[leftmargin=1.4em, itemsep=0.35em]
  \item PI/lead --- MSCA CASREx fellowship project
  \item Programme lead --- TransGP fabric manipulation (2024--2025)
  \item Assistant Supervisor --- 2 MEng placement students, Imperial (2026)
  \item Mentoring --- UG FYP and PhD students at CUHK; GTA at Auckland
  \item Workshop organiser \& presenter --- IEEE ICMA 2025
\end{itemize}

\subsection*{\sffamily\color{accent} Impact --- academic}
Invited talks (Imperial Dept.\ Seminar, Robotics PG Committee, WRC 2025); IEEE reviewer; RiTA best paper (2017); ETA 2026 selection; GATE 99.12\%.

\subsection*{\sffamily\color{accent} Impact --- social \& economic}
Stent testing for implant manufacturers (BvW); garment manufacturing automation (sewing, alignment, flattening); upper-arm rehabilitation exosuits (CASREx); medical simulators (RoSE, SoGut).

\subsection*{\sffamily\color{accent} Collaborations \& network}
Imperial (Morph Lab); TransGP/HKU (Prof.\ Kosuge); CUHK (Prof.\ Lau); Auckland (Prof.\ Xu, Prof.\ Cheng); LS2N Nantes (Prof.\ Caro); multi-country track: India $\to$ NZ $\to$ Hong Kong $\to$ China $\to$ UK.

\subsection*{\sffamily\color{accent} Teaching \& supervision}
Guest lecture --- PyDrake simulation \& control, DESE71002, Imperial (2026); MAEG5090 Topics in Robotics, CUHK (50 students); 4 Auckland UG courses as GTA; 9 UG + 1 PG courses as Lecturer, Galgotias.

\vspace{0.6em}
\section*{Elevator top 5 --- for a 2-minute panel}

\begin{enumerate}[leftmargin=1.6em, itemsep=0.65em, label=\rank{\arabic*.}]

  \item \textbf{I lead a Horizon Europe Marie Sk\l{}odowska-Curie Fellowship (CASREx, EUR~276k) at Imperial,}
  developing cable-driven exosuits for upper-arm muscle-impairment rehabilitation with human trials, EMG, and digital twins.

  \item \textbf{I have a sustained publication record in top robotics and control venues:}
  \textit{Soft Robotics}, \textit{IEEE TIE}, \textit{IEEE TCST}, \textit{MMT}, \textit{IEEE RA-L}, \textit{IEEE T-ASE}, and \textit{IROS} --- spanning healthcare, cable robots, and garment manufacturing.

  \item \textbf{My work translates beyond academia:}
  industry-funded esophageal stent testing (BvW, NZD~39k $\to$ RoSE platform), and leading fabric-robotics R\&D at TransGP for real garment-manufacturing workflows.

  \item \textbf{I lead research, not only contribute:}
  MSCA project PI; led the TransGP fabric programme; now Assistant Supervisor for two MEng students; mentor and grant co-author from CUHK onwards.

  \item \textbf{I am gaining European and international recognition:}
  European Talent Academy 2026; invited speaker at World Robotics Conference 2025; active reviewer for RA-L, ICRA, and JMR.

\end{enumerate}

\vspace{0.5em}
\begin{center}
  \small\color{charcoal}
  Source: \texttt{cv-simple.pdf} and section files in \texttt{CV\_latex\_source/} \textbar{}
  \href{https://profiles.imperial.ac.uk/d.bhattacharya1/about}{Imperial Profile}
\end{center}

\end{document}
